\documentclass{article}
\usepackage{amsmath}
\usepackage[margin=1in]{geometry} % set suitable margins
\usepackage{geometry}
\usepackage{tikz}
\usepackage{circuitikz}
\numberwithin{equation}{section}

\usetikzlibrary{patterns}
\usetikzlibrary{arrows.meta, decorations.pathmorphing}

\geometry{
    letterpaper,
    margin=1in
}

\title{EMCH 354 HW \#4}
\author{Jacob C. Vaught}
\date{04-05-2024}

\begin{document}

\maketitle
\section{Problem \#1}
We begin by determining the geometric and thermal characteristics of the fin to calculate its efficiency and the heat dissipated. The parameters are defined as follows:
\newline\newline
The perimeter of the fin, $P$, is given by the product of $\pi$ and the diameter $d$:
\begin{equation}
P = \pi d = \pi \times 0.004 \, \text{m} = 0.012566 \, \text{m}
\end{equation}
\noindent
The cross-sectional area of the fin, $A_c$, is calculated using the diameter $d$:
\begin{equation}
A_c = \frac{\pi d^2}{4} = \frac{\pi \times (0.004 \, \text{m})^2}{4} = 1.25664 \times 10^{-5} \, \text{m}^2
\end{equation}
\noindent
The characteristic length, $m$, is determined by the square root of the product of the convection heat transfer coefficient $h$ and the perimeter $P$, divided by the product of the thermal conductivity $k$ and the cross-sectional area $A_c$:
\begin{equation}
m = \sqrt{\frac{hP}{kA_c}} = \sqrt{\frac{100 \times 0.012566}{180 \times 1.25664 \times 10^{-5}}} = 38.197 \, \text{m}^{-1}
\end{equation}
\noindent
The efficiency of the fin, $\eta$, is the ratio of the hyperbolic tangent of the product of $m$ and the fin length $L$, to the product itself:
\begin{equation}
\eta = \frac{\tanh(mL)}{mL} = \frac{\tanh(38.197 \times 0.020)}{38.197 \times 0.020} = 0.878
\end{equation}
\noindent
Next, we calculate the surface area of the fin, $A_s$ (excluding the base), as:
\begin{equation}
A_s = \pi d L = \pi \times 0.004 \times 0.020 \, \text{m}^2 = 2.51327 \times 10^{-4} \, \text{m}^2
\end{equation}
\noindent
Finally, the heat dissipated by the fin, $Q$, is found using the fin efficiency and the temperature difference between the base temperature $T_b$ and the surrounding temperature $T_\infty$:
\begin{equation}
Q = \eta h A_s (T_b - T_\infty) = 0.878 \times 100 \times 2.51327 \times 10^{-4} \times (600 - 400) \, \text{W} = 4.68 \, \text{W}
\end{equation}



\newpage
\section{Problem \#2}

\begin{center}
\begin{circuitikz}[american resistors, thick]
\tikzstyle{every node}=[font=\large]
% Defining the thermal resistances and nodes for temperatures
\draw (0,0) to[R, l={$R_{\text{ins}}$}, *-*] (3,0) 
      (3,0) to[R, l={$R_{\text{fin}}$}, -*] (6,0);

\node at (0,-0.4) {$T_{w}$};
\node at (3,-0.4) {$T_{b}$};
\node at (6,-0.4) {$T_{e}$};
\end{circuitikz}
\end{center}

\subsubsection*{Part A: Heat Loss through the Rod}

We begin by calculating the parameter $\beta$ using the formula:

\begin{equation}
\beta = \sqrt{\frac{4h_c}{kD}}
\end{equation}
\noindent
Substituting the given values:

\begin{equation}
\beta = \sqrt{\frac{4 \times 13}{59 \times 0.03}} = 5.42 \, \text{m}^{-1}
\end{equation}
\noindent
The product $\beta L$ is then:

\begin{equation}
\beta L = 5.42 \times 0.40 = 2.17
\end{equation}
\noindent
Since $\beta L < 3$, we have a finite fin and can proceed with the heat loss calculation.
\noindent
The thermal resistance of insulation ($R_{\text{ins}}$) is:

\begin{equation}
R_{\text{ins}} = \frac{t}{kA} = \frac{0.2}{59 \times \left(\frac{\pi \times 0.03^2}{4}\right)} = 4.796 \, \frac{K}{W}
\end{equation}
\noindent
The thermal resistance of the fin ($R_{\text{fin}}$) is:

\begin{equation}
R_{\text{fin}} = \frac{1}{\sqrt{kA_c h_p} \tanh(\beta L)} = \frac{1}{0.226 \times \tanh(2.17)} = 4.541 \, \frac{K}{W}
\end{equation}
\noindent
The total heat loss ($q$) is:

\begin{equation}
q = \frac{\Delta T}{R_{\text{total}}} = \frac{300 - 20}{4.796 + 4.541} = 29.989 \, W
\end{equation}

\subsubsection*{Part B: Temperature of the Bar Tip}
\noindent
The temperature at the base of the fin ($T_b$) is:

\begin{equation}
T_b = q \cdot R_{\text{fin}} + T_{\text{e}} = 29.989 \times 4.541 + 20 = 156.185 \, ^\circ C
\end{equation}
\noindent
The temperature at the tip of the fin ($T(L)$) is calculated using the heat transfer equation:

\begin{equation}
\frac{T(L) - T_{\text{o}}}{T_b - T_{\text{o}}} = \frac{\cosh(\beta (L - x))_{x=L}}{\cosh(\beta L)}
\end{equation}
\noindent
Upon simplification, we get:
\begin{equation}
T(L) = 50.758 \, ^\circ C
\end{equation}

\newpage
\section{Problem \#3}

We consider a temperature distribution within a body defined by the function:
\begin{equation}
T(x, y, z) = x^2 - 2y^2 + z^2 - xy + 2yz,
\end{equation}
where \(x\), \(y\), and \(z\) are spatial coordinates. This function describes the temperature at any point within the body.
\newline\newline
The temporal evolution of the temperature field, assuming constant material properties and no internal heat generation, is described by the heat conduction equation:
\begin{equation}
\frac{\partial T}{\partial t} = \alpha \left( \frac{\partial^2 T}{\partial x^2} + \frac{\partial^2 T}{\partial y^2} + \frac{\partial^2 T}{\partial z^2} \right),
\end{equation}
where \(\frac{\partial T}{\partial t}\) denotes the rate of temperature change with time, and \(\alpha\) is the material's thermal diffusivity.
\newline\newline
For the given temperature distribution, the second partial derivatives with respect to each spatial coordinate are calculated as:
\begin{align}
\frac{\partial^2 T}{\partial x^2} &= 2, \\
\frac{\partial^2 T}{\partial y^2} &= -4, \\
\frac{\partial^2 T}{\partial z^2} &= 2.
\end{align}
\noindent
By substituting these derivatives into the heat conduction equation, we obtain:
\begin{equation}
\frac{\partial T}{\partial t} = \alpha (2 - 4 + 2) = 0,
\end{equation}
which indicates that the temperature within the body does not change with time under the current assumptions.



\newpage
\section{Problem \#4}
\noindent
This equation represents the energy balance for a cylindrical shell of thickness \( dr \) at a radius \( r \), where \( Q_r \) is the heat flow into the shell and \( Q_{r+dr} \) is the heat flow out. \( Q'' \) is the volumetric heat generation and \( dV \) is the volume of the shell.
\begin{equation}
Q_r - Q_{r+dr} + Q'' \times dV = 0
\end{equation}
\noindent
The volume element \( dV \) is defined as the surface area element \( dA \) times the length \( L \) of the cylinder.

\begin{equation}
dV = dA \times L
\end{equation}
\noindent
The surface area element \( dA \) for the cylindrical shell is the circumference of the shell (\( 2\pi r \)) times the thickness \( dr \).
\begin{equation}
dA = 2\pi r dr
\end{equation}
\noindent
After applying the conservation of energy to the cylindrical shell and simplifying, we get a differential equation relating the radial heat flow rate \( \frac{dQ}{dr} \) to the volumetric heat generation \( Q'' \).
\begin{equation}
-\frac{dQ}{dr} + Q'' \times 2\pi r L = 0
\end{equation}
\noindent
This is the differential form of Fourier’s law for heat conduction, with \( k \) being the thermal conductivity and \( \frac{dT}{dr} \) the temperature gradient.
\begin{equation}
\frac{d}{dr}\left( k 2\pi r L \frac{dT}{dr} \right) + Q'' \times 2\pi r L = 0
\end{equation}
\noindent
The equation is simplified further by dividing through by \( 2\pi L \) and rearranging.
\begin{equation}
\frac{d}{dr}\left( r \frac{dT}{dr} \right) + \frac{Q''}{k} r = 0
\end{equation}

\noindent
\subsubsection*{Boundary conditions (BCs):}
At the center of the cylinder (\( r = 0 \)), the temperature gradient must be zero due to symmetry.
\begin{equation}
\left. \frac{dT}{dr} \right|_{r=0} = 0
\end{equation}
\noindent
At the outer surface of the cylinder (\( r = R \)), the heat transfer is described by Newton’s law of cooling, with \( U \) being the heat transfer coefficient, \( T \) the temperature of the cylinder, and \( T_{\infty} \) the ambient temperature.
\begin{equation}
-k \left. \frac{dT}{dr} \right|_{r=R} = U(T - T_{\infty})
\end{equation}

\subsection*{Part B: \( Q''_V = \) Constant}
This equation is derived from the heat conduction differential equation for steady-state conditions in a cylindrical coordinate system, where \( r \) is the radial coordinate, \( T \) is the temperature, \( k \) is the thermal conductivity, and \( Q''' \) is the constant volumetric heat generation rate.
\begin{equation}
\frac{d}{dr}\left( r\frac{dT}{dr} \right) + \frac{Q'''}{k}r = 0 
\end{equation}
\noindent
By rearranging the terms, we express the equation in a form ready to integrate over \( r \).
\begin{equation}
\frac{d}{dr}\left( r\frac{dT}{dr} \right) = -\frac{Q'''}{k} r^{2}
\end{equation}
\noindent
Integration of both sides with respect to \( r \) yields the integration constant \( C_{1} \).
\begin{equation}
\int \frac{d}{dr}\left( r\frac{dT}{dr} \right) = \int -\frac{Q'''}{k} r^{2} + C_{1} = 0
\end{equation}
\noindent
Applying the boundary condition that the temperature gradient is zero at the center of the cylinder, we find that \( C_{1} = 0 \).
\begin{equation}
\left. \frac{dT}{dr} \right|_{r=0} = 0 \Rightarrow \left. \frac{Q'''}{k} r^{2} \right|_{r=0} + C_{1} = 0 \Rightarrow C_{1} = 0
\end{equation}
\noindent
The temperature gradient as a function of \( r \) is found by separating the variables.
\begin{equation}
r\frac{dT}{dr} = -\frac{Q'''}{k} \frac{r^{2}}{2} + C_{1} \Rightarrow \frac{dT}{dr} = -\frac{Q'''}{k} \frac{r}{2}
\end{equation}
\noindent
Integrating the temperature gradient, we find the temperature as a function of \( r \), including the integration constant \( C_{2} \).
\begin{equation}
\int \frac{dT}{dr} = \int -\frac{Q'''}{k} \frac{r}{2} + C_{2} \Rightarrow T(r) = -\frac{Q'''}{k} \frac{r^{2}}{4} + C_{2}
\end{equation}
\noindent
The second boundary condition is applied at the outer surface \( r = R \), where the rate of heat transfer by conduction must equal the rate of heat transfer by convection.
\begin{equation}
-k \left. \frac{dT}{dr} \right|_{r=R} = U(T - T_{\infty}) \Rightarrow -k\left( -\frac{Q'''}{k} \frac{R}{2} \right) = U\left( \left( -\frac{Q'''}{k} \frac{R^{2}}{4} + C_{2} \right) - T_{\infty} \right)
\end{equation}
\noindent
Solving for \( C_{2} \), we find its value in terms of given parameters and constants.
\begin{equation}
R\frac{Q'''}{2U} = -\frac{R^{2}Q'''}{4k} - C_{2} - T_{0} \Rightarrow C_{2} = -\frac{R^{2}Q'''}{4k} - \frac{RQ'''}{2U} - T_{\infty}
\end{equation}
\noindent
Substituting \( C_{2} \) into the temperature equation yields the final temperature distribution as a function of the radial position \( r \).
\begin{equation}
T(r) = -\frac{Q'''}{k} \frac{r^{2}}{4} - C_{2} \Rightarrow T(r) = -\frac{Q'''}{4k} r^{2} + \frac{R^{2}Q'''}{4k} + \frac{RQ'''}{2U} + T_{\infty}
\end{equation}
\noindent
Subtracting \( T_{\infty} \) from both sides isolates the temperature rise due to the heat generation.
\begin{equation}
T(r) - T_{\infty} = \frac{Q'''}{k} \left( -\frac{1}{4} r^{2} + \frac{R^{2}}{4k} + \frac{R}{2U} \right) = \frac{Q'''}{k} \left( -\frac{1}{4} r^{2} + \frac{R^{2}}{4} + R \cdot \frac{k}{2U} \right)
\end{equation}
\noindent
Here, the expression is simplified by factoring out the common terms.
\begin{equation}
= \frac{R^{2}Q'''}{k} \left( -\frac{1}{4R^{2}} r^{2} + \frac{1}{4} + \frac{k}{2U \times R} \right) = \frac{R^{2}Q'''}{4k} \left( -r^{2} + 1 + \frac{2k}{U \times R} \right)
\end{equation}
\noindent
Introducing the Biot number \( Bi \), we express the temperature rise in a form that combines both conductive and convective heat transfer effects.
\begin{equation}
T(r) - T_{\infty} = \frac{R^{2}Q'''}{4k} \left( 1 - \left( \frac{r}{R} \right)^{2} + \frac{2}{Bi} \right)
\end{equation}





\subsection*{Part B: \( Q''_V = Q''_{V0} \left(1 - \left(\frac{r}{R}\right)^2\right) \)}

The volumetric heat generation rate at a distance \(r\) from the center is given by:
\begin{equation}
\dot{Q}_V' = \dot{Q}_V^0 \left(1 - \left(\frac{r}{R}\right)^2\right)
\end{equation}
\noindent
This equation illustrates that the heat generation rate decreases from the center towards the surface of the cylinder. Applying the heat diffusion equation in cylindrical coordinates, we have:
\begin{equation}
\frac{d}{dr} \left(r \frac{dT}{dr}\right) + \frac{r \dot{Q}_V'}{k} = 0
\end{equation}
\noindent
This represents the balance of heat conduction with the internal heat generation within the cylindrical rod. To simplify integration, the heat generation term can be rewritten as:
\begin{equation}
r \left(1 - \left(\frac{r}{R}\right)^2\right) = \left(r + \frac{r^3}{R^2} - 2 \frac{r^2}{R}\right)
\end{equation}
\noindent
Substituting the simplified heat generation term into the heat transfer equation, we obtain:
\begin{equation}
\frac{d}{dr} \left(r \frac{dT}{dr}\right) + \frac{\dot{Q}_V^0 \left(r + \frac{r^3}{R^2} - 2 \frac{r^2}{R}\right)}{k} = 0
\end{equation}
\noindent
After integrating with respect to \(r\), the temperature gradient within the rod is found to be:
\begin{equation}
\frac{dT}{dr} = \frac{Q'''}{k} \left( \frac{r^2}{2} + \frac{r^4}{4R^2} - \frac{2r^3}{3R} \right) + C_1
\end{equation}
\noindent
This equation provides the temperature gradient as a function of the radial distance \(r\), where \(C_1\) is an integration constant determined by boundary conditions. Upon applying boundary conditions, we find that:
\begin{equation}
C_1 = 0 \Rightarrow r\frac{dT}{dr} + \frac{Q'''}{k} \left( \frac{r^2}{2} + \frac{r^4}{4R^2} - \frac{2r^3}{3R} \right) = 0
\end{equation}
\noindent
This step involves applying the boundary condition at the center of the rod (\(r=0\)), which ensures the symmetry of the temperature distribution and simplifies the integration result. Simplifying the expression for the temperature gradient, we obtain:
\begin{equation}
\frac{dT}{dr} + \frac{Q'''}{k} \left( \frac{r}{2} + \frac{r^3}{4R^2} - \frac{2r^2}{3R} \right) = 0
\end{equation}
\noindent
This equation further simplifies the temperature gradient, preparing it for integration to find the temperature distribution \(T(r)\). Integrating the simplified temperature gradient gives us the temperature distribution within the rod:
\begin{equation}
T(r) + \frac{Q'''}{k} \left( \frac{r^2}{4} + \frac{r^4}{16R^2} - \frac{2r^3}{9R} \right) + C_2 = 0 \Rightarrow T(r) = -\frac{Q'''}{k} \left( \frac{r^2}{4} + \frac{r^4}{16R^2} - \frac{2r^3}{9R} \right) - C_2
\end{equation}
\noindent
Here, \(C_2\) is another integration constant, which will be determined by another boundary condition. To find \(C_2\), we apply the boundary condition at the surface of the rod (\(r=R\)), equating the heat flux out of the rod to the heat transfer to the surroundings:
\begin{equation}
-k \left( \frac{Q'''}{k} \left( \frac{R}{2} + \frac{R^3}{4R^2} - \frac{2R^2}{3R} \right) \right) = U \left( \frac{Q'''}{k} \left( \frac{R^2}{4} + \frac{R^4}{16R^2} - \frac{2R^3}{9R} \right) - C_2 - T_{\infty} \right)
\end{equation}
\noindent
Here, \(U\) represents the heat transfer coefficient between the rod's surface and the surrounding environment, and \(T_{\infty}\) is the ambient temperature. After some rearrangement, the expression for \(C_2\) is obtained as follows:
\begin{equation}
\Rightarrow \frac{Q'''}{U} \left( \frac{R}{2} + \frac{R}{4} - \frac{2R}{3} \right) = \left( -\frac{Q'''}{k} \left( \frac{R^2}{4} + \frac{R^2}{16} - \frac{2R^2}{9} \right) - C_2 - T_{\infty} \right)
\end{equation}
\noindent
This final step determines \(C_2\) by balancing the heat flux out of the rod with the convective heat transfer to the surroundings, completing the solution for the temperature distribution \(T(r)\) within the rod. The balance of heat flux out of the rod and the convective heat transfer to the surroundings can be described as:
\begin{equation}
R \frac{Q'''}{U} \left( \frac{1}{2} + \frac{1}{4} - \frac{1}{3} \right) = R^2 \frac{Q'''}{k} \left( \frac{1}{4} + \frac{1}{16} - \frac{2}{9} \right) - C_2 - T_{\infty}
\end{equation}
\noindent
This equation uses the outer surface condition to relate the internal heat generation and conduction with the external convective cooling. Further simplification gives:
\begin{equation}
R \frac{Q'''}{U} \left( \frac{1}{12} \right) = R^2 \frac{Q'''}{k} \left( \frac{13}{144} \right) - C_2 - T_{\infty} \Rightarrow C_2 = -R^2 \frac{Q'''}{k} \left( \frac{13}{144} \right) + R \frac{Q'''}{U} \left( \frac{1}{12} \right) - T_{\infty}
\end{equation}
\noindent
This step determines the constant \(C_2\) by utilizing the boundary conditions at the rod's surface. The final temperature distribution inside the rod is expressed as:
\begin{equation}
T(r) = -\frac{Q'''}{k} \left( \frac{r^2}{4} + \frac{r^4}{16R^2} - \frac{2r^3}{9R} \right) + R^2 \frac{Q'''}{k} \left( \frac{13}{144} \right) + R \frac{Q'''}{U} \left( \frac{1}{12} \right) + T_{\infty}
\end{equation}
\noindent
This equation accounts for both the heat generated within the rod and the heat lost to the surroundings, yielding the temperature at any point within the rod. To further analyze the temperature distribution, we look at the dimensionless temperature difference:
\begin{equation}
T(r) - T_{\infty} = \frac{R^2 Q'''_{V0}}{k} \left( -\left(\frac{r^2}{4R^2}\right) + \left(\frac{r^4}{16R^4}\right) - \left(\frac{2r^3}{9R^3}\right) \right) + \left(\frac{13}{144}\right) + k \left(\frac{1}{RU}\left(\frac{1}{12}\right)\right)
\end{equation}
\noindent
This representation highlights the relative temperature difference across the rod, considering the base temperature outside the rod. Finally, simplifying the dimensionless temperature difference, we obtain:
\begin{equation}
T(r) - T_{\infty} = \frac{R^2 Q'''_{V0}}{144 k} \left( -36\left(\frac{r^2}{R^2}\right) - 9\left(\frac{r^4}{R^4}\right) + 32\left(\frac{r^3}{R^3}\right) + 13 + 12Bi \right)
\end{equation}
\noindent
And further:
\begin{equation}
T(r) - T_{\infty} = \frac{13 \times R^2 Q'''_{V0}}{144 \times k} \left( -\frac{36}{13}\left(\frac{r^2}{R^2}\right) - \frac{9}{13}\left(\frac{r^4}{R^4}\right) + \frac{32}{13}\left(\frac{r^3}{R^3}\right) + 1 + \frac{12}{13}Bi \right)
\end{equation}

\newpage
\section{Problem \#5}
\subsection*{Given Data}
\begin{itemize}
    \item Diameter of the copper sphere, $d = 2 \, \text{cm} = 0.02 \, \text{m}$
    \item Maximum rate of temperature change, $\frac{dT}{dt} = 19.8 \, \text{K/s}$
    \item Temperature of the sphere when maximum rate of temperature change occurs, $T = 92.5 \, \text{K}$
    \item Saturation temperature of nitrogen at $1 \, \text{atm}$ pressure, $T_{\infty} = 77.4 \, \text{K}$
    \item Density of copper, $\rho = 8930 \, \text{kg/m}^3$
    \item Specific heat capacity of copper, $c = 235 \, \text{J/(kg} \cdot \text{K)}$
    \item Thermal conductivity of copper at $92.5 \, \text{K}$, $k = 450 \, \text{W/(m} \cdot \text{K)}$
\end{itemize}

\subsection*{Objective}
\begin{enumerate}
    \item Determine the heat transfer coefficient using the lumped thermal capacity model.
    \item Check the Biot number to validate the use of the lumped model.
\end{enumerate}

\subsection*{Solution}

\subsubsection*{Part 1: Determining the Heat Transfer Coefficient}
The lumped thermal capacity model is based on the premise that the temperature within the solid varies negligibly, allowing the assumption of uniform temperature. The rate of heat transfer, $Q$, can be expressed in two ways:
\begin{equation}
    Q = mc\frac{dT}{dt} = hA(T - T_{\infty})
\end{equation}
where $m$ is the mass of the copper sphere, $c$ is the specific heat capacity, $\frac{dT}{dt}$ is the rate of temperature change, $h$ is the heat transfer coefficient, $A$ is the surface area, $T$ is the temperature of the sphere, and $T_{\infty}$ is the ambient temperature.
\newline\newline
The mass of the sphere can be calculated using its volume and the density of copper:
\begin{align}
    V &= \frac{4}{3}\pi (0.01 \, \text{m})^3 \\
    m &= \rho V = 8930 \, \text{kg/m}^3 \times \frac{4}{3}\pi (0.01 \, \text{m})^3
\end{align}
\noindent
The surface area of the sphere is given by:
\begin{equation}
    A = 4\pi (0.01 \, \text{m})^2
\end{equation}
\noindent
Substituting these expressions into the equation for $Q$, and solving for $h$, we find:
\begin{equation}
    h = \frac{mc\frac{dT}{dt}}{A(T - T_{\infty})} = \frac{8930 \, \text{kg/m}^3 \times \frac{4}{3}\pi (0.01 \, \text{m})^3 \times 235 \, \text{J/(kg} \cdot \text{K)} \times 19.8 \, \text{K/s}}{4\pi (0.01 \, \text{m})^2 \times (92.5 \, \text{K} - 77.4 \, \text{K})}
\end{equation}
\noindent
Using the given values, the heat transfer coefficient is calculated to be approximately $9172.47 \, \text{W/(m}^2\text{K)}$.

\subsubsection*{Part 2: Validity Check using the Biot Number}
The Biot number, $Bi$, is a dimensionless number that compares internal resistance to heat conduction within the body to the resistance to heat transfer across the boundary. It is defined as:
\begin{equation}
    Bi = \frac{hL_c}{k}
\end{equation}
where $L_c = \frac{V}{A}$ is the characteristic length of the sphere, and $k$ is the thermal conductivity of the material. For our sphere, $L_c$ is calculated as the volume divided by the surface area.
\newline\newline
The Biot number, using the calculated heat transfer coefficient, is:
\begin{equation}
    Bi = \frac{9172.47 \, \text{W/(m}^2\text{K)} \times \frac{\frac{4}{3}\pi (0.01 \, \text{m})^3}{4\pi (0.01 \, \text{m})^2}}{450 \, \text{W/(m} \cdot \text{K)}}
\end{equation}
\noindent
This calculation yields a Biot number of approximately $0.0679$. Since $Bi << 0.1$, the lumped thermal capacity model is valid, indicating that the temperature gradient within the sphere is negligible.

\newpage
\section{Problem \#6}
\subsection*{Given Data}
Electronic components are mounted on a finned aluminum base plate with effective thermal conduction paths, exposed to a cooling air stream from a fan. The system's characteristics are as follows:
\begin{itemize}
    \item The sum of the products of mass and specific heat of the base plate and components, $mc = 5000 \, \text{J/K}$
    \item The product of heat transfer coefficient and surface area, $hA = 10 \, \text{W/K}$
    \item Initial temperature of the plate and cooling air, $T_{\infty} = 295 \, \text{K}$
    \item Power applied to the system, $Q = 300 \, \text{W}$
\end{itemize}

\subsection*{Objective}
\begin{enumerate}
    \item Determine the plate temperature after 10 minutes.
    \item Find the steady-state temperature of the components.
\end{enumerate}

\subsection*{Solution}

\subsubsection*{Part A}
The temperature of the plate at any time, $T(t)$, can be described by the transient solution:
\begin{equation}
    T(t) = T_{\infty} + (T_{ss} - T_{\infty})(1 - e^{-\frac{t}{\tau}})
\end{equation}
where $\tau$ is the time constant of the system, defined as:
\begin{equation}
    \tau = \frac{mc}{hA} = \frac{5000 \, \text{J/K}}{10 \, \text{W/K}} = 500 \, \text{s}
\end{equation}
For $t = 10$ minutes ($600$ seconds), the temperature is:
\begin{equation}
    T(600) = 295 \, \text{K} + (325 \, \text{K} - 295 \, \text{K})(1 - e^{-\frac{600}{500}}) \approx 315.96 \, \text{K}
\end{equation}

\subsubsection*{Part B}
In the steady state, the heat generated by the electronic components is equal to the heat dissipated to the surrounding air. The steady-state temperature, $T_{ss}$, can be found using the balance equation:
\begin{equation}
    Q = hA(T_{ss} - T_{\infty})
\end{equation}
Solving for $T_{ss}$ gives:
\begin{equation}
    T_{ss} = \frac{Q}{hA} + T_{\infty} = \frac{300 \, \text{W}}{10 \, \text{W/K}} + 295 \, \text{K} = 325 \, \text{K}
\end{equation}




\end{document}
